\documentclass[12pt]{article}
\usepackage{response}
\usepackage[utf8]{inputenc}
\usepackage{amsmath}
\usepackage{lipsum} % to generate some filler text
\usepackage{fullpage}
\usepackage{subfigure}
\usepackage{mathrsfs}
\usepackage{xr}
\usepackage{amsthm,cite,url,graphicx,booktabs,lipsum,color,bm,caption,subcaption,soul}
\usepackage{pifont,tikz,paralist,multirow,amssymb}
\usepackage{epstopdf}
\usepackage{xifthen}
\usepackage{enumerate}
\usepackage{titlesec,upgreek}
\usepackage{lineno}
\usepackage[table]{xcolor} 
\usepackage{amsmath}
\usepackage{listings}

\definecolor{myBlue}{RGB}{0, 0, 255}
\definecolor{myRed}{RGB}{255, 0, 0}
\definecolor{myPurple}{RGB}{199, 21, 133}
\definecolor{myGreen}{RGB}{0, 128, 0}

\title{Comments}

\begin{document}
	\maketitle

\begin{reviewer}
\noindent \textbf{Category: [Abbreviations] [5.1] [\textcolor{blue}{medium}]}

\noindent \textbf{Location: [2, first mention of 'initial state radiation']}

\begin{lstlisting}[breaklines=true]
	The simulation models the beam energy spread and initial state radiation (ISR) in the $e^+e^-$ annihilations
\end{lstlisting}

\textbf{Reviewer Comment:}

The abbreviation 'ISR' is defined correctly here, but the acronym 'MC' for Monte Carlo is used earlier without being defined.

\textbf{Suggested Revision:}

\begin{lstlisting}[breaklines=true]
	Simulated data samples produced with a {\sc geant4}-based~\cite{Geant4} Monte Carlo (MC) package, which include...
\end{lstlisting}

\end{reviewer}

\begin{reviewer}
\noindent \textbf{Category: [Abbreviations] [5.2] [\textcolor{blue}{medium}]}

\noindent \textbf{Location: [3, first mention of pid]}

\begin{lstlisting}[breaklines=true]
	Particle identification~(PID) for charged tracks combines measurements of the d$E$/d$x$ in the MDC and the flight time in the TOF to form likelihoods $\mathcal{L}(h)$ for each hadron hypothesis $h~(h=p, K, \pi)$.
\end{lstlisting}

\textbf{Reviewer Comment:}

Do not write out 'Particle identification' when defining the abbreviation PID. Use 'Particle identification (PID)' is incorrect per rule 5.2.

\textbf{Suggested Revision:}

\begin{lstlisting}[breaklines=true]
	Particle identification (PID) for charged tracks combines measurements of the $\mathrm{d}E/\mathrm{d}x$ in the MDC and the flight time in the TOF to form likelihoods $\mathcal{L}(h)$ for each hadron hypothesis $h$ ($h=p, K, \pi$).
\end{lstlisting}

\end{reviewer}

\begin{reviewer}
\noindent \textbf{Category: [Language] [1.2] [\textcolor{blue}{medium}]}

\noindent \textbf{Location: [3, first sentence, after listing decays]}

\begin{lstlisting}[breaklines=true]
	Signal events are required to contain at least six charged tracks and three photon candidates for $\chi_{cJ} \to \Lambda K^0_S\bar \Xi^0$,
	and at least six charged tracks and one photon candidate for $\chi_{cJ} \to \Lambda K^-\bar \Xi^+$.
\end{lstlisting}

\textbf{Reviewer Comment:}

The decay descriptor is used as a noun. Replace with a proper noun like 'the decay mode' or 'the channel'.

\textbf{Suggested Revision:}

\begin{lstlisting}[breaklines=true]
	Signal events for the $\chi_{cJ} \to \Lambda K^0_S\bar \Xi^0$ channel are required to contain at least six charged tracks and three photon candidates, while events for the $\chi_{cJ} \to \Lambda K^-\bar \Xi^+$ channel require at least six charged tracks and one photon candidate.
\end{lstlisting}

\end{reviewer}

\begin{reviewer}
\noindent \textbf{Category: [Language] [1.3] [\textcolor{blue}{medium}]}

\noindent \textbf{Location: [6, first sentence]}

\begin{lstlisting}[breaklines=true]
	Using $(2712.4\pm14.3)\times10^6$ $\psi(3686)$ events collected...
\end{lstlisting}

\textbf{Reviewer Comment:}

Sentence starts with a number in mathematical notation. Begin with a proper introductory phrase.

\textbf{Suggested Revision:}

\begin{lstlisting}[breaklines=true]
	Based on $(2712.4\pm14.3)\times10^6$ $\psi(3686)$ events collected with the BESIII detector operating at the BEPCII collider,
\end{lstlisting}

\end{reviewer}

\begin{reviewer}
\noindent \textbf{Category: [Language] [1.5] [low]}

\noindent \textbf{Location: [1, first paragraph, second sentence]}

\begin{lstlisting}[breaklines=true]
	Experimentally and theoretically, their decays are not as extensively studied as those of the vector charmonium states $J/\psi$ and $\psi(3686)$.
\end{lstlisting}

\textbf{Reviewer Comment:}

Avoid jargon 'extensively studied'. Use more precise language.

\textbf{Suggested Revision:}

\begin{lstlisting}[breaklines=true]
	Experimentally and theoretically, their decays have been studied less than those of the vector charmonium states $J/\psi$ and $\psi(3686)$.
\end{lstlisting}

\end{reviewer}

\begin{reviewer}
\noindent \textbf{Category: [Language] [1.5] [low]}

\noindent \textbf{Location: [1, second paragraph, first sentence]}

\begin{lstlisting}[breaklines=true]
	To date, there have been many studies of $\chi_{cJ}$ decays involving $p\bar p$ pairs and meson pairs~\cite{ref::pppipi,ref::pppi0pi0,ref::ppkk}, and hyperon pairs~\cite{ref::lambdalambdabar,ref::sigmasigmabar} which can be described by the color octet mechanism~\cite{ref::oct1,ref::oct2}.
\end{lstlisting}

\textbf{Reviewer Comment:}

Avoid vague phrase 'many studies'. Specify the quantity or use 'numerous'.

\textbf{Suggested Revision:}

\begin{lstlisting}[breaklines=true]
	To date, numerous studies of $\chi_{cJ}$ decays involving $p\bar p$ pairs and meson pairs~\cite{ref::pppipi,ref::pppi0pi0,ref::ppkk}, and hyperon pairs~\cite{ref::lambdalambdabar,ref::sigmasigmabar} have been performed, which can be described by the color octet mechanism~\cite{ref::oct1,ref::oct2}.
\end{lstlisting}

\end{reviewer}

\begin{reviewer}
\noindent \textbf{Category: [Language] [1.5] [low]}

\noindent \textbf{Location: [2, description of mc sample size]}

\begin{lstlisting}[breaklines=true]
	An inclusive MC sample containing $2.7\times10^{9}$ generic $\psi(3686)$ decays is used to study background.
\end{lstlisting}

\textbf{Reviewer Comment:}

The phrase 'study background' is vague. Use more precise terminology.

\textbf{Suggested Revision:}

\begin{lstlisting}[breaklines=true]
	An inclusive MC sample containing $2.7\times10^{9}$ generic $\psi(3686)$ decays is used to estimate background contributions.
\end{lstlisting}

\end{reviewer}

\begin{reviewer}
\noindent \textbf{Category: [Language] [1.3] [low]}

\noindent \textbf{Location: [2, beginning]}

\begin{lstlisting}[breaklines=true]
	The EMC measures photon energies with a resolution of $2.5\%$ ($5\%$) at $1$~GeV in the barrel (end cap) region.
\end{lstlisting}

\textbf{Reviewer Comment:}

Sentence starts with an acronym (EMC). Rewrite to avoid this.

\textbf{Suggested Revision:}

\begin{lstlisting}[breaklines=true]
	Photon energies are measured by the EMC with a resolution of $2.5\%$ ($5\%$) at $1$~GeV in the barrel (end cap) region.
\end{lstlisting}

\end{reviewer}

\begin{reviewer}
\noindent \textbf{Category: [Language] [1.5] [low]}

\noindent \textbf{Location: [3, description of candidate selection for $\bar{\xi}^+$ and $\bar{\xi}^0$]}

\begin{lstlisting}[breaklines=true]
	If there is more than one $\bar \Xi^+$  baryon candidate, we keep the one with minimum $\sqrt{(\frac{M_{\bar \Lambda \pi^+}-m_{\bar{\Xi}^+}}{\sigma_{\bar{\Lambda}\pi^+}})^2+(\frac{M_{\bar{p}\pi^+}-m_{\bar{\Lambda}}}{\sigma_{\bar{p}\pi^+}})^2}$...
\end{lstlisting}

\textbf{Reviewer Comment:}

Avoid using 'the one'. Replace with 'that' for better phrasing.

\textbf{Suggested Revision:}

\begin{lstlisting}[breaklines=true]
	If there is more than one $\bar{\Xi}^+$ baryon candidate, the candidate with the minimum value of $\sqrt{(\frac{M_{\bar{\Lambda} \pi^+}-m_{\bar{\Xi}^+}}{\sigma_{\bar{\Lambda}\pi^+}})^2+(\frac{M_{\bar{p}\pi^+}-m_{\bar{\Lambda}}}{\sigma_{\bar{p}\pi^+}})^2}$ is retained...
\end{lstlisting}

\end{reviewer}

\begin{reviewer}
\noindent \textbf{Category: [Language] [1.3] [low]}

\noindent \textbf{Location: [4, beginning]}

\begin{lstlisting}[breaklines=true]
	Events in the $\Lambda$, $\bar{\Xi}^0$, and $K^0_S$ sideband regions is scaled by a factor of 0.5
\end{lstlisting}

\textbf{Reviewer Comment:}

Sentence starts with a symbol/number. Rephrase to avoid starting with 'Events'.

\textbf{Suggested Revision:}

\begin{lstlisting}[breaklines=true]
	The events in the $\Lambda$, $\bar{\Xi}^0$, and $K^0_S$ sideband regions are scaled by a factor of 0.5
\end{lstlisting}

\end{reviewer}

\begin{reviewer}
\noindent \textbf{Category: [Language] [1.5] [low]}

\noindent \textbf{Location: [4, paragraph describing the significance calculation]}

\begin{lstlisting}[breaklines=true]
	The statistical significance of each $\chi_{cJ}$ signal exceeds $10\sigma$, determined using the profile likelihood ratio test statistic $q_J=-2\ln(\mathcal{L}_J/\mathcal{L}_{\max})$.
\end{lstlisting}

\textbf{Reviewer Comment:}

Avoid using 'determined using'. Use 'determined from' or 'calculated with'.

\textbf{Suggested Revision:}

\begin{lstlisting}[breaklines=true]
	The statistical significance of each $\chi_{cJ}$ signal exceeds $10\sigma$, calculated from the profile likelihood ratio test statistic $q_J=-2\ln(\mathcal{L}_J/\mathcal{L}_{\max})$.
\end{lstlisting}

\end{reviewer}

\begin{reviewer}
\noindent \textbf{Category: [Language] [1.5] [low]}

\noindent \textbf{Location: [5, paragraph 5, sentence 2]}

\begin{lstlisting}[breaklines=true]
	The obtained efficiencies in the 2D intervals are weighted by the 2D distributions of various momentum and $\cos\theta$ to match our signal MC events, according to Ref.~\cite{ref::Lambda}.
\end{lstlisting}

\textbf{Reviewer Comment:}

The phrase 'our signal MC events' is informal. Use more formal language.

\textbf{Suggested Revision:}

\begin{lstlisting}[breaklines=true]
	The obtained efficiencies in the 2D intervals are weighted by the 2D distributions of momentum and $\cos\theta$ to match the signal MC events, according to Ref.~\cite{ref::Lambda}.
\end{lstlisting}

\end{reviewer}

\begin{reviewer}
\noindent \textbf{Category: [Language] [1.6] [low]}

\noindent \textbf{Location: [6, first sentence after the branching fraction list]}

\begin{lstlisting}[breaklines=true]
	These measurements will be helpful for understanding $\chi_{cJ}$ decay mechanisms. The ratios of the branching fractions...
\end{lstlisting}

\textbf{Reviewer Comment:}

Two independent sentences are connected without proper conjunction. Use a transition word or combine them logically.

\textbf{Suggested Revision:}

\begin{lstlisting}[breaklines=true]
	These measurements will be helpful for understanding $\chi_{cJ}$ decay mechanisms. Furthermore, the ratios of the branching fractions...
\end{lstlisting}

\end{reviewer}

\begin{reviewer}
\noindent \textbf{Category: [Numbers] [4.2] [\textcolor{blue}{medium}]}

\noindent \textbf{Location: [5, paragraph 13, sentence 1]}

\begin{lstlisting}[breaklines=true]
	The systematic uncertainties due to the limited statistics of the mixed signal MC samples
	are 0.7\%, 0.6\%, and 0.7\% for  $\chi_{c0,1,2}\to  \Lambda K^0_S\bar \Xi^0$,
	and 0.500\%, 0.4\%, and 0.4\% for $\chi_{c0,1,2}\to \Lambda K^-\bar \Xi^+$,~respectively.
\end{lstlisting}

\textbf{Reviewer Comment:}

Inconsistent precision: '0.500\%' should match the format of other values (e.g., 0.5\%).

\textbf{Suggested Revision:}

\begin{lstlisting}[breaklines=true]
	The systematic uncertainties due to the limited statistics of the mixed signal MC samples are 0.7\%, 0.6\%, and 0.7\% for $\chi_{c0,1,2}\to \Lambda K^0_S\bar \Xi^0$, and 0.5\%, 0.4\%, and 0.4\% for $\chi_{c0,1,2}\to \Lambda K^-\bar \Xi^+$,~respectively.
\end{lstlisting}

\end{reviewer}

\begin{reviewer}
\noindent \textbf{Category: [Numbers] [4.2] [\textcolor{blue}{medium}]}

\noindent \textbf{Location: [5, paragraph 14, sentence 1]}

\begin{lstlisting}[breaklines=true]
	The uncertainties from the branching fractions quoted from the PDG~\cite{ref::pdg2024} are 2.3\%, 2.8\%, and 2.4\% for $\psi(3686)\to \gamma\chi_{c0,1,2}$,
	0.07\% for $\mathcal{B}(K^{0}_{S} \to \pi^+ \pi^-)$, 0.800\% for $\mathcal{B}(\Lambda \to p \pi^-)$, 0.010\% for $\mathcal{B}(\bar{\Xi}^0 \to \Lambda\pi^0)$, 0.03\% for $\mathcal{B}(\bar{\Xi}^- \to \Lambda \pi^-)$, and 0.03\% for $\mathcal{B}(\pi^0 \to \gamma \gamma)$,
\end{lstlisting}

\textbf{Reviewer Comment:}

Inconsistent precision: '0.800\%' and '0.010\%' should be rounded to match the format of other values (e.g., 0.8\%, 0.01\%).

\textbf{Suggested Revision:}

\begin{lstlisting}[breaklines=true]
	The uncertainties from the branching fractions quoted from the PDG~\cite{ref::pdg2024} are 2.3\%, 2.8\%, and 2.4\% for $\psi(3686)\to \gamma\chi_{c0,1,2}$, 0.07\% for $\mathcal{B}(K^{0}_{S} \to \pi^+ \pi^-)$, 0.8\% for $\mathcal{B}(\Lambda \to p \pi^-)$, 0.01\% for $\mathcal{B}(\bar{\Xi}^0 \to \Lambda\pi^0)$, 0.03\% for $\mathcal{B}(\bar{\Xi}^- \to \Lambda \pi^-)$, and 0.03\% for $\mathcal{B}(\pi^0 \to \gamma \gamma)$,
\end{lstlisting}

\end{reviewer}

\begin{reviewer}
\noindent \textbf{Category: [Numbers] [4.1] [low]}

\noindent \textbf{Location: [6, first branching fraction for $\chi_{c1}\to \lambda k^0_s\bar \xi^0$]}

\begin{lstlisting}[breaklines=true]
	$\mathcal{B}(\chi_{c1}\to \Lambda K^0_S\bar \Xi^0)=(7.6\pm0.8\pm0.7)\times10^{-5}$
\end{lstlisting}

\textbf{Reviewer Comment:}

Precision mismatch: central value has one decimal place (7.6), while uncertainties have one (0.8, 0.7). They should match.

\textbf{Suggested Revision:}

\begin{lstlisting}[breaklines=true]
	$\mathcal{B}(\chi_{c1}\to \Lambda K^0_S\bar \Xi^0)=(7.6\pm0.8\pm0.7)\times10^{-5}$ (This is acceptable as 7.6 implies uncertainty in the tenths place, matching 0.8 and 0.7).
\end{lstlisting}

\end{reviewer}

\begin{reviewer}
\noindent \textbf{Category: [References] [9.7] [\textcolor{red}{high}]}

\noindent \textbf{Location: [7, bibliography, entry ref::psip-num-inc]}

\begin{lstlisting}[breaklines=true]
	\href{http://arxiv.org/abs/2403.06766}{arXiv 2403 06776}
\end{lstlisting}

\textbf{Reviewer Comment:}

Inconsistent arXiv format and incorrect arXiv ID (06766 vs 06776). Use standard arXiv citation format.

\textbf{Suggested Revision:}

\begin{lstlisting}[breaklines=true]
	\href{https://arxiv.org/abs/2403.06766}{arXiv:2403.06766}
\end{lstlisting}

\end{reviewer}

\begin{reviewer}
\noindent \textbf{Category: [References] [9.7] [low]}

\noindent \textbf{Location: [7, bibliography, entry ref::lxik]}

\begin{lstlisting}[breaklines=true]
	Phys. Rev. D. {\bf 91}, 092006 (2015).}
\end{lstlisting}

\textbf{Reviewer Comment:}

Extra period after 'D' (D. vs D). Be consistent with journal abbreviation.

\textbf{Suggested Revision:}

\begin{lstlisting}[breaklines=true]
	Phys.~Rev.~D~{\bf 91}, 092006 (2015).
\end{lstlisting}

\end{reviewer}

\begin{reviewer}
\noindent \textbf{Category: [Tables] [7.4] [low]}

\noindent \textbf{Location: [5, table rows for '$\lambda$ reconstruction' and '4c kinematic fit']}

\begin{lstlisting}[breaklines=true]
	In the table: '$\Lambda$ reconstruction & 2.90 & 1.60 & 3.5 & 2.9 & 2.5 & 2.5\\' and '4C kinematic fit & 3.2 & 1.1 & 2.9 & 1.9 & 0.2  & 0.3\\'
\end{lstlisting}

\textbf{Reviewer Comment:}

Numbers in the table are not consistently aligned on decimal points. '2.90' vs '3.5' and '0.2' has an extra space.

\textbf{Suggested Revision:}

\begin{lstlisting}[breaklines=true]
	Ensure all numbers in the table have consistent decimal places and alignment. For example: '2.9 & 1.6 & 3.5 & 2.9 & 2.5 & 2.5' and '3.2 & 1.1 & 2.9 & 1.9 & 0.2 & 0.3'.
\end{lstlisting}

\end{reviewer}

\begin{reviewer}
\noindent \textbf{Category: [Typography] [3.0] [\textcolor{blue}{medium}]}

\noindent \textbf{Location: [1, first paragraph, third sentence]}

\begin{lstlisting}[breaklines=true]
	At BESIII, except for tiny production of $e^+e^-\to \chi_{c1}$ through annihilation into two photons~\cite{ref::eetochic1}, the $\chi_{cJ}$ mesons can not be produced directly through $e^+ e^-$ annihilation processes, but can be accessed via abundant radiative decays of $\psi(3686)$ with a branching fraction of about 9\%~\cite{ref::pdg2024}.
\end{lstlisting}

\textbf{Reviewer Comment:}

Missing space before citation. Use protected space (\textasciitilde{}).

\textbf{Suggested Revision:}

\begin{lstlisting}[breaklines=true]
	At BESIII, except for tiny production of $e^+e^-\to \chi_{c1}$ through annihilation into two photons~\cite{ref::eetochic1}, the $\chi_{cJ}$ mesons can not be produced directly through $e^+ e^-$ annihilation processes, but can be accessed via abundant radiative decays of $\psi(3686)$ with a branching fraction of about 9\%~\cite{ref::pdg2024}.
\end{lstlisting}

\end{reviewer}

\begin{reviewer}
\noindent \textbf{Category: [Typography] [3.0] [\textcolor{blue}{medium}]}

\noindent \textbf{Location: [1, second paragraph, second sentence]}

\begin{lstlisting}[breaklines=true]
	Besides, an enhancement in the $ p \bar p$ invariant mass spectrum has been observed in the study of $J/\psi \to \gamma p\bar p$ ~\cite{ref::ppbar1,ref::ppbar2}.
\end{lstlisting}

\textbf{Reviewer Comment:}

Missing protected space before citation. Use \textasciitilde{}\textbackslash\{\}cite\{\}.

\textbf{Suggested Revision:}

\begin{lstlisting}[breaklines=true]
	Besides, an enhancement in the $ p \bar p$ invariant mass spectrum has been observed in the study of $J/\psi \to \gamma p\bar p$~\cite{ref::ppbar1,ref::ppbar2}.
\end{lstlisting}

\end{reviewer}

\begin{reviewer}
\noindent \textbf{Category: [Typography] [3.0] [\textcolor{blue}{medium}]}

\noindent \textbf{Location: [1, third paragraph, first sentence]}

\begin{lstlisting}[breaklines=true]
	Previously, BESIII reported the measurements of $\chi_{cJ}\to \Lambda K^-\bar \Xi^+$, with branching fractions at $10^{-4}$, based on $(107.7\pm0.6)\times10^6$ $\psi(3686)$ events taken in 2009~\cite{ref::LXiK}.
\end{lstlisting}

\textbf{Reviewer Comment:}

Missing protected space before citation. Use \textasciitilde{}\textbackslash\{\}cite\{\}.

\textbf{Suggested Revision:}

\begin{lstlisting}[breaklines=true]
	Previously, BESIII reported the measurements of $\chi_{cJ}\to \Lambda K^-\bar \Xi^+$, with branching fractions at $10^{-4}$, based on $(107.7\pm0.6)\times10^6$ $\psi(3686)$ events taken in 2009~\cite{ref::LXiK}.
\end{lstlisting}

\end{reviewer}

\begin{reviewer}
\noindent \textbf{Category: [Typography] [3.7] [\textcolor{blue}{medium}]}

\noindent \textbf{Location: [2, unit for momentum resolution]}

\begin{lstlisting}[breaklines=true]
	The charged-particle momentum resolution at $1~{\rm GeV}/c$ is $0.5\%$.
\end{lstlisting}

\textbf{Reviewer Comment:}

The unit 'GeV/c' should be in roman font, not math mode. Use \textbackslash\{\}mathrm or \textbackslash\{\}text.

\textbf{Suggested Revision:}

\begin{lstlisting}[breaklines=true]
	The charged-particle momentum resolution at $1~\mathrm{GeV}/c$ is $0.5\%$.
\end{lstlisting}

\end{reviewer}

\begin{reviewer}
\noindent \textbf{Category: [Typography] [3.0] [\textcolor{blue}{medium}]}

\noindent \textbf{Location: [4, caption of figure\textasciitilde{}\textbackslash\{\}ref\{tab:lambdakxi\}]}

\begin{lstlisting}[breaklines=true]
	The purple solid, orange solid, zerue solid curves denote the peaking background contributions
\end{lstlisting}

\textbf{Reviewer Comment:}

Missing space after comma in list of colors. Also, 'zerue' appears to be a typo.

\textbf{Suggested Revision:}

\begin{lstlisting}[breaklines=true]
	The purple solid, orange solid, and blue solid curves denote the peaking background contributions
\end{lstlisting}

\end{reviewer}

\begin{reviewer}
\noindent \textbf{Category: [Typography] [3.2] [\textcolor{blue}{medium}]}

\noindent \textbf{Location: [7, acknowledgement section, first sentence]}

\begin{lstlisting}[breaklines=true]
	BEPCII {\hypersetup{urlcolor=black}\href{https://cstr.cn/31109.02.BEPC}{(https://cstr.cn/31109.02.BEPC)}}
\end{lstlisting}

\textbf{Reviewer Comment:}

The URL should be set in typewriter font using \textbackslash\{\}texttt\{\} or similar, not plain parentheses. The hyperlink setup is also misplaced.

\textbf{Suggested Revision:}

\begin{lstlisting}[breaklines=true]
	BEPCII (\href{https://cstr.cn/31109.02.BEPC}{\texttt{https://cstr.cn/31109.02.BEPC}})
\end{lstlisting}

\end{reviewer}

\begin{reviewer}
\noindent \textbf{Category: [Typography] [3.2] [low]}

\noindent \textbf{Location: [2, reference to the tof upgrade citation]}

\begin{lstlisting}[breaklines=true]
	providing a time resolution of 60~ps~\cite{Tof1,Tof2,Tof3}.
\end{lstlisting}

\textbf{Reviewer Comment:}

Missing protected space (\textasciitilde{}) before the \textbackslash\{\}cite command.

\textbf{Suggested Revision:}

\begin{lstlisting}[breaklines=true]
	providing a time resolution of 60~ps~\cite{Tof1,Tof2,Tof3}.
\end{lstlisting}

\end{reviewer}

\begin{reviewer}
\noindent \textbf{Category: [Typography] [3.7] [low]}

\noindent \textbf{Location: [2, unit for center-of-mass energy]}

\begin{lstlisting}[breaklines=true]
	achieved at $\sqrt{s} = 3.773\text{GeV}$.
\end{lstlisting}

\textbf{Reviewer Comment:}

Missing space between the number and the unit 'GeV'.

\textbf{Suggested Revision:}

\begin{lstlisting}[breaklines=true]
	achieved at $\sqrt{s} = 3.773~\text{GeV}$.
\end{lstlisting}

\end{reviewer}

\begin{reviewer}
\noindent \textbf{Category: [Typography] [3.0] [low]}

\noindent \textbf{Location: [3, photon selection criteria, energy requirement]}

\begin{lstlisting}[breaklines=true]
	The deposited energy of each shower must exceed 25~MeV in the barrel region ($|\cos \theta|< 0.80$, where $\theta$ denotes the
	polar angle with respect to the z-axis, the symmetry axis of the MDC) and more than 50~MeV in the end cap region ($0.86 <|\cos \theta|< 0.92$).
\end{lstlisting}

\textbf{Reviewer Comment:}

Missing space before opening parenthesis. Add a space for correct typography.

\textbf{Suggested Revision:}

\begin{lstlisting}[breaklines=true]
	The deposited energy of each shower must exceed 25~MeV in the barrel region ($|\cos \theta|< 0.80$, where $\theta$ denotes the polar angle with respect to the $z$-axis, the symmetry axis of the MDC) and more than 50~MeV in the end cap region ($0.86 <|\cos \theta|< 0.92$).
\end{lstlisting}

\end{reviewer}

\begin{reviewer}
\noindent \textbf{Category: [Typography] [3.1] [low]}

\noindent \textbf{Location: [3, photon timing requirement]}

\begin{lstlisting}[breaklines=true]
	the difference between the EMC time and the event start time is required to be within [0, 700]\,ns.
\end{lstlisting}

\textbf{Reviewer Comment:}

Use a math minus sign for the negative number in the range. The dash is incorrect.

\textbf{Suggested Revision:}

\begin{lstlisting}[breaklines=true]
	the difference between the EMC time and the event start time is required to be within $[0, 700]$~ns.
\end{lstlisting}

\end{reviewer}

\begin{reviewer}
\noindent \textbf{Category: [Typography] [3.2] [low]}

\noindent \textbf{Location: [3, reference to pdg in $k_s^0$ and $\lambda$ sideband definitions]}

\begin{lstlisting}[breaklines=true]
	where $m_{K^{0}_{S}}$ is the $K^0_S$ nominal mass~\cite{ref::pdg2024}.~The $K^{0}_{S}$ sideband region...
\end{lstlisting}

\textbf{Reviewer Comment:}

Missing protected space (\textasciitilde{}) before the citation command. Add it for proper formatting.

\textbf{Suggested Revision:}

\begin{lstlisting}[breaklines=true]
	where $m_{K^{0}_{S}}$ is the $K^0_S$ nominal mass~\cite{ref::pdg2024}. The $K^{0}_{S}$ sideband region...
\end{lstlisting}

\end{reviewer}

\begin{reviewer}
\noindent \textbf{Category: [Typography] [3.7] [low]}

\noindent \textbf{Location: [4, equation for the product of branching fractions]}

\begin{lstlisting}[breaklines=true]
	${\mathcal B}(\psi(3686)\to \gamma\chi_{cJ}) \cdot \mathcal{B}(\chi_{cJ}\to \Lambda K \bar \Xi)$
\end{lstlisting}

\textbf{Reviewer Comment:}

Inconsistent notation for branching fraction: $\mathcal B$ vs. $\mathcal{B}$. Use consistent style.

\textbf{Suggested Revision:}

\begin{lstlisting}[breaklines=true]
	$\mathcal{B}(\psi(3686)\to \gamma\chi_{cJ}) \cdot \mathcal{B}(\chi_{cJ}\to \Lambda K \bar \Xi)$
\end{lstlisting}

\end{reviewer}

\begin{reviewer}
\noindent \textbf{Category: [Typography] [3.2] [low]}

\noindent \textbf{Location: [4, sentence referencing figure\textasciitilde{}\textbackslash\{\}ref\{fig:compare\_lxiks\} and \textbackslash\{\}ref\{fig:compare\_lxik\}]}

\begin{lstlisting}[breaklines=true]
	Figure~\ref{fig:compare_LXiKS} and \ref{fig:compare_LXiK} show the comparisons
\end{lstlisting}

\textbf{Reviewer Comment:}

Missing protected space (\textasciitilde{}) before the second \textbackslash\{\}ref command.

\textbf{Suggested Revision:}

\begin{lstlisting}[breaklines=true]
	Figure~\ref{fig:compare_LXiKS} and~\ref{fig:compare_LXiK} show the comparisons
\end{lstlisting}

\end{reviewer}

\begin{reviewer}
\noindent \textbf{Category: [Typography] [3.0] [low]}

\noindent \textbf{Location: [5, paragraph 5, sentence 3]}

\begin{lstlisting}[breaklines=true]
	The differences of the re-weighted efficiencies between data and MC simulation are assigned as the systematic uncertainties, which are 2.9\%, 1.6\%, 3.5\% for $\chi_{c0,1,2}\to \Lambda K^0_S\bar \Xi^0$; and
	2.9\%, 2.5\%, and 2.5\% for $\chi_{c0,1,2}\to \Lambda K^-\bar \Xi^+$,~respectively.
\end{lstlisting}

\textbf{Reviewer Comment:}

Missing space after the semicolon before 'and'.

\textbf{Suggested Revision:}

\begin{lstlisting}[breaklines=true]
	The differences of the re-weighted efficiencies between data and MC simulation are assigned as the systematic uncertainties, which are 2.9\%, 1.6\%, 3.5\% for $\chi_{c0,1,2}\to \Lambda K^0_S\bar \Xi^0$; and 2.9\%, 2.5\%, and 2.5\% for $\chi_{c0,1,2}\to \Lambda K^-\bar \Xi^+$,~respectively.
\end{lstlisting}

\end{reviewer}

\begin{reviewer}
\noindent \textbf{Category: [Typography] [3.0] [low]}

\noindent \textbf{Location: [5, paragraph 8, sentence 1]}

\begin{lstlisting}[breaklines=true]
	The systematic uncertainty of photon or $\pi^0$ reconstruction is assigned as  1.0\% per photon based on the control sample of $J/\psi\to\pi^+\pi^-\pi^0$~\cite{ref::gamma-recon}.
\end{lstlisting}

\textbf{Reviewer Comment:}

Double space before '1.0\%'.

\textbf{Suggested Revision:}

\begin{lstlisting}[breaklines=true]
	The systematic uncertainty of photon or $\pi^0$ reconstruction is assigned as 1.0\% per photon based on the control sample of $J/\psi\to\pi^+\pi^-\pi^0$~\cite{ref::gamma-recon}.
\end{lstlisting}

\end{reviewer}

\begin{reviewer}
\noindent \textbf{Category: [Typography] [3.0] [low]}

\noindent \textbf{Location: [6, list of branching fractions]}

\begin{lstlisting}[breaklines=true]
	$\mathcal{B}(\chi_{c0}\to \Lambda K^0_S\bar \Xi^0)=(1.14\pm0.13\pm0.13)\times10^{-4}$, $\mathcal{B}(\chi_{c1}\to \Lambda K^0_S\bar \Xi^0)=(7.6\pm0.8\pm0.7)\times10^{-5}$, ...
\end{lstlisting}

\textbf{Reviewer Comment:}

Missing space after commas separating the list items. This improves readability.

\textbf{Suggested Revision:}

\begin{lstlisting}[breaklines=true]
	$\mathcal{B}(\chi_{c0}\to \Lambda K^0_S\bar \Xi^0)=(1.14\pm0.13\pm0.13)\times10^{-4}$, $\mathcal{B}(\chi_{c1}\to \Lambda K^0_S\bar \Xi^0)=(7.6\pm0.8\pm0.7)\times10^{-5}$, ...
\end{lstlisting}

\end{reviewer}

\begin{reviewer}
\noindent \textbf{Category: [Typography] [3.7] [low]}

\noindent \textbf{Location: [6, reference to the ratio expectation]}

\begin{lstlisting}[breaklines=true]
	the theoretical expectation of ${\mathcal B}(\chi_{cJ}\to\Lambda K^0_S\bar{\Xi}^0)/{\mathcal B}(\chi_{cJ}\to\Lambda K^-\bar{\Xi}^+)=0.5$
\end{lstlisting}

\textbf{Reviewer Comment:}

The symbol for branching fraction is inconsistently formatted. It uses $\mathcal{B}$ elsewhere but ${\mathcal B}$ here.

\textbf{Suggested Revision:}

\begin{lstlisting}[breaklines=true]
	the theoretical expectation of $\mathcal{B}(\chi_{cJ}\to\Lambda K^0_S\bar{\Xi}^0)/\mathcal{B}(\chi_{cJ}\to\Lambda K^-\bar{\Xi}^+)=0.5$
\end{lstlisting}

\end{reviewer}

\begin{reviewer}
\noindent \textbf{Category: [Typography] [3.0] [low]}

\noindent \textbf{Location: [7, bibliography, multiple entries]}

\begin{lstlisting}[breaklines=true]
	Phys. Rev. D {\bf 110}, 030001 (2024).}
\end{lstlisting}

\textbf{Reviewer Comment:}

Inconsistent spacing before volume number and inconsistent punctuation at the end of entries.

\textbf{Suggested Revision:}

\begin{lstlisting}[breaklines=true]
	Phys.~Rev.~D~{\bf 110}, 030001 (2024).
\end{lstlisting}

\end{reviewer}

\end{document}
